
\def\PUDcodigo{EME107}

\if\COMPILAR0
   \def\PUDcodacad{04507.16}
\else
   \def\PUDcodacad{04506.14}
\fi

\def\PUDdisciplina{Metrologia}

\def\PUDsemestre{S3}

\def\PUDhoras{80}

%NOVOS ---------------------------------------------------
\def\PUDhorasTeoria{50}
\def\PUDhorasPratica{30}

\def\PUDinterECA{---}

\def\PUDatuaisECA{---}

\def\PUDrecursos{Projetor multimídia; Quadro branco e pincel; Aulas em laboratório de metrologia.}

%----------------------------------------------------------

\def\PUDcreditos{4}

\def\PUDprerequisito{---}

\def\PUDpreacad{---}

\def\PUDobjetivos{Conhecer as definições e terminologias da metrologia;  Compreender e avaliar os parâmetros envolvidos em um processo de medição;  Conhecer as principais técnicas e instrumentos/sistemas de medição;  Compreender a importância da metrologia nos processos industriais.}

\def\PUDementa{Terminologias e Conceitos da Metrologia;  Unidades de Medidas e Sistema Internacional;  Instrumentos Convencionais de Medição;  Medição Ótica 2D;  Medição Tridimensional por Coordenadas;  Engenharia Reversa Através de Scanner a Laser;  Sistemas de Medição;  Sistemas de Tolerâncias e Ajustes;  Tolerâncias Geométricas;  Calibração;  Incerteza da Medição.}

\def\PUDconteudo{ & Linguagem, conceitos e terminologias da metrologia.  Erros de medição.  Resultado de uma medição. 
\\ 
 & Unidades do sistema internacional de unidades.  Grafia das unidades.  Unidades derivadas.  Fatores de conversão.  Constantes.  Dimensão de uma grandeza.  
\\ 
 & Medições com o uso de paquímetros, micrômetros, relógio comparador, projetor de perfil e goniômetro. 
\\ 
 & Visualização ótica bidimensional para reprodução de contronos e medição em componentes de pequeno porte. 
\\ 
 & Medição tridimensional por coordenadas para controle de qualidade em componentes de médio a grande porte.  Engenharia reversa, por meio de scanner a laser, em peças de médio a grande porte. 
\\ 
 & Métodos básicos de medição.  Módulos de um sistema de medição.  Características metrológicas dos sistemas de medição.  Representação absoluta e relativa. 
\\ 
 & Tolerâncias.  Afastamentos inferior e superior.  Dimensões mínimas e máximas.  Folgas mínimas e máximas.  Interferências mínimas e máximas. 
\\ 
 & Superfícies.  Tolerância geométrica de forma.  Tolerância geométrica de posição.  Tolerância de batidura. 
\\ 
 & Procedimentos de medição.  Registros de medições.  Certificados de calibração.  
\\ 
 & Conceitos básicos em estatística.  Propagação de erros.  Incerteza da medição.  Incerteza da medição expandida.  Modelo matemático.  Fatores que contribuem para a incerteza da medição. }

\def\PUDmetodo{Aulas expositórias que podem ser teóricas e/ou práticas, onde as práticas no laboratório serão marcadas no decorrer da disciplina.}

\def\PUDavalia{A avaliação será feita com aplicação de provas teóricas e/ou práticas, além da possibilidade de inclusão de trabalhos e seminários no decorrer da disciplina.}

\def\PUDbibbasica{ & \href{www.biblioteca.ifce.edu.br/index.asp?codigo_sophia=24518}{Albertazzi}, A.; Sousa, A. R.  Fundamentos de Metrologia Científica e Industrial.  1ª ed.  Barueri, SP: Manole, 2008.  408p.  ISBN: 9788520421161. 
\\ 
 & \href{www.biblioteca.ifce.edu.br/index.asp?codigo_sophia=65387}{Lira}, F. A.  Metrologia na Indústria.  10ª ed.  São Paulo, SP: Érica, 2015.  256p.  ISBN: 9788571947832. 
\\ 
 &  \href{www.biblioteca.ifce.edu.br/index.asp?codigo_sophia=62371}{Silva Neto}, J. C.  Metrologia e Controle Dimensional - Conceitos, Normas e Aplicações.  1ª ed.  Rio de Janeiro, RJ: Elsevier, 2012.  239p.  ISBN: 9788535255799. }

\def\PUDbibcomplementar{ &  \href{www.biblioteca.ifce.edu.br/index.asp?codigo_sophia=24552}{Bini}, Edson.  A Técnica da ajustagem: metrologia, medição, roscas, acabamento.  São Paulo, SP: Hemus, 2004.  210p.  ISBN: 8528905284. 
\\ 
 & \href{www.biblioteca.ifce.edu.br/index.asp?codigo_sophia=24606}{Balbinot}, A.; Brusamarello, V. J.  Instrumentação e fundamentos de medidas 1: princípios e definições.  Rio de Janeiro, RJ: LTC, 2006.  477p.  ISBN: 8521614969. 
\\ 
 &  \href{www.biblioteca.ifce.edu.br/index.asp?codigo_sophia=24137}{Balbinot}, A.; Brusamarello, V. J.  Instrumentação e fundamentos de medidas 2 : medição de pressão.  Rio de Janeiro, RJ: LTC, 2007.  658p.  ISBN: 9788521615637.
\\
 &  \href{www.biblioteca.ifce.edu.br/index.asp?codigo_sophia=57333}{Toledo}, José Carlos de.  Sistemas de medição e metrologia.  E-book.   Intersaberes.  196p.  ISBN: 9788582129418.
\\
 &  \href{www.biblioteca.ifce.edu.br/index.asp?codigo_sophia=58862}{Santos}, Josiane Oliveira dos.  Metrologia e normalização.  Pearson.  124p.  ISBN: 9788543016757. }
%\\ 
 %&  GUEDES, P.;  Metrologia Industrial.  1ª edição.  Rio de Janeiro.  Editora ETEP BRASIL.  2011.  424p.  ISBN 9789728480271. }

\def\PUDautor{Francisco Nélio Costa Freitas}

\def\PUDdata{2013-05-22}

\def\PUDrevisao{1}

\def\PUDrevisor{Samuel Vieira Dias}

\def\PUDdatarevisao{2017-05-11}

\PUDcorpo
